\subsection{Task 6 - Recognize characters}
\subsubsection{Recognition}

The final stage of the system consists of identifying each segmented and preprocessed character by comparing it against a set of reference templates (alphabet) obtained in Task 5.

\subsubsection{Methodology}

Each input character image is first normalized to a fixed size of $32 \times 32$ pixels using the preprocessing pipeline described in Task 4. This ensures consistency between the input data and the stored templates.

The recognition process is based on template matching using a similarity metric. For each candidate character in the alphabet, a similarity score is computed between the normalized input image $X$ and the corresponding template $T_i$. In this implementation, the similarity is measured using the normalized two-dimensional correlation:

\begin{equation}
S_i = \text{corr2}(X, T_i)
\end{equation}

where $S_i \in [-1, 1]$, with higher values indicating greater similarity.

The system evaluates all available templates in the alphabet structure, which is implemented as a MATLAB struct where each field corresponds to a character label (e.g., \texttt{alphabet.A}, \texttt{alphabet.B}, etc.).

The predicted character $\hat{c}$ is selected as the one that maximizes the similarity score:

\begin{equation}
\hat{c} = \arg\max_i S_i
\end{equation}

\subsubsection{Confidence Criterion}

To improve robustness and avoid incorrect classifications, a confidence threshold is introduced. If the maximum similarity score is below a predefined threshold (empirically set to 0.5), the recognition is considered unreliable and the output is rejected:

\begin{equation}
\text{if } \max(S_i) < \tau \Rightarrow \hat{c} = \text{``?''}
\end{equation}

This mechanism prevents low-confidence matches from being accepted as valid predictions.

Additionally, a margin-based criterion comparing the best and second-best scores can be introduced to further improve classification reliability.

\subsubsection{Discussion}

The use of normalized correlation provides a simple and computationally efficient approach to character recognition under controlled conditions. Since all characters share the same font and size, template matching is sufficient to achieve high accuracy.

However, this method presents several limitations. First, it is highly sensitive to segmentation errors and misalignments, as even small shifts can significantly reduce the correlation score. Second, the use of a single template per character does not capture intra-class variability, which may lead to incorrect recognition in the presence of noise or distortions.

To address these limitations, more advanced approaches could be considered, such as storing multiple templates per character, using normalized feature descriptors, or applying machine learning-based classifiers.

Despite these limitations, the implemented approach provides a solid baseline solution for the recognition task within the constraints of the problem.
